\section{Motivation}
\label{sec:motivation}

% 2-3 pages.

\subsection{Anatomy of the convergence loop}
% Walk a real ticket: edit, lint, typecheck, unit, integration, review.
% 50-200 edit-run cycles. Emphasise: the tree mutates between every gate run.
% Name the two kinds of step the loop contains and keep them apart from here
% on: ACTIONS change the tree (the agent's edit, a formatter's fix); GATES judge
% it (lint, typecheck, tests). An action is done or not done; a gate's verdict
% is about one snapshot and no other.

\subsection{Why one-shot execution fails}
% The core observation. A DAG engine node fires once when dependencies are
% satisfied. A gate must fire repeatedly as the tree changes.
% NOTE: do NOT frame this as "graphs vs DAGs" -- our dependency graph IS acyclic
% and admission-enforced to be so (CEL rules plus controller-side plan
% resolution; there are no admission webhooks). The contrast is EXECUTION
% SEMANTICS:
%   one-shot topological execution  vs.  content-addressed level-triggered convergence.
% Table: topology / trigger / node identity / completion criterion.

\subsection{Why agent frameworks fail}
% No kernel isolation, no capacity model, no multi-tenancy, no declarative
% workflow surface, no tracker-neutral integration.

\subsection{Requirements}
% R1 verification results must be re-entrant and content-keyed
% R2 invalidation must be path-scoped or cost grows superlinearly -- and it
%    must fail closed, or scoping becomes a way to skip verification
% R3 agent-authored artifacts must not run on shared kernels
% R4 the agent must not be able to weaken its own verification (this becomes:
%    policy-owned gates, class-only behaviour, spec the agent cannot write)
% R5 capacity is finite and non-elastic at the isolation layer
% R6 tracker integration must not leak into the core
% R7 the decision procedure must be pure: testable without a cluster and
%    replayable after a controller crash
